% ============================================================
% Chapter 8 — System Integration, Verification, Validation & Deployment
% Source: PSYCON Engineering Design Document, Vol. 1, Ch. 8
% ============================================================
\section{System Integration, Verification, Validation, and Deployment}
\label{sec:integration-verification}

\subsection{System Integration Strategy}
\label{sec:integration-strategy}

Integration occurs in five stages, illustrated in \Cref{fig:integration-stages}. Every stage must pass before proceeding to the next.

\begin{figure}[H]
\centering
\begin{tikzpicture}[
    font=\scriptsize,
    node distance=0mm,
    field/.style={draw, rounded corners=2pt, fill=blue!7, minimum width=3.6cm, minimum height=0.65cm, align=center, line width=0.5pt}
]
\node[field] (s1) {Individual Sensor Validation};
\node[field, below=6mm of s1] (s2) {Subsystem Integration};
\node[field, below=6mm of s2] (s3) {Complete Hardware Assembly};
\node[field, below=6mm of s3] (s4) {Backend Integration};
\node[field, below=6mm of s4] (s5) {End-to-End Validation};

\draw[-{Stealth[length=2mm]}, line width=0.6pt] (s1.south) -- (s2.north);
\draw[-{Stealth[length=2mm]}, line width=0.6pt] (s2.south) -- (s3.north);
\draw[-{Stealth[length=2mm]}, line width=0.6pt] (s3.south) -- (s4.north);
\draw[-{Stealth[length=2mm]}, line width=0.6pt] (s4.south) -- (s5.north);
\end{tikzpicture}
\caption{Five-stage PSYCON integration strategy, from individual sensor validation through to end-to-end validation.}
\label{fig:integration-stages}
\end{figure}

\subsection{Stage 1 --- Individual Sensor Validation}
\label{sec:stage1-sensor-validation}

Each module is tested independently before integration. \Cref{tab:sensor-validation} summarizes the verification checks and acceptance criteria for each sensor.

\begin{table*}[t]
\centering
\caption{Stage 1 individual sensor validation: verification checks and acceptance criteria.}
\label{tab:sensor-validation}
\small
\begin{tabularx}{\textwidth}{@{}l Y Y@{}}
\toprule
\textbf{Sensor} & \textbf{Verification Checks} & \textbf{Acceptance Criteria} \\
\midrule
MAX30102 & I\textsuperscript{2}C communication; raw RED and IR readings; finger detection; stable waveform; no communication errors & Device detected at \texttt{0x57}; stable signal; no bus failures \\
MPU6050 & Accelerometer; gyroscope; temperature register; axis orientation & Address \texttt{0x68}; stable readings at rest; motion reflected correctly \\
MCP9808 & Ambient temperature; stable readings; repeatability & Address \texttt{0x18}; no sudden fluctuations \\
ADS1115 & Address detection; all four channels readable; stable ADC values & Address \texttt{0x48}; no random saturation \\
TEMT6000 & Analog response to light; dark vs. bright environments & Smooth output variation \\
INMP441 & Audio recording; no clipping; correct I\textsuperscript{2}S configuration; acceptable noise floor & Clean recordings; no DMA overruns \\
\bottomrule
\end{tabularx}
\end{table*}

\subsection{Stage 2 --- I\texorpdfstring{\textsuperscript{2}}{2}C Bus Validation}
\label{sec:stage2-i2c-validation}

An I\textsuperscript{2}C scanner is run across the Wrist Module bus. The expected devices and addresses match those previously listed in \Cref{tab:i2c-addresses}.

\begin{checklist}
    \item All devices detected.
    \item No duplicate addresses.
    \item No intermittent failures.
\end{checklist}

\subsection{Stage 3 --- Integrated Wrist Module}
\label{sec:stage3-wrist-integration}

All wrist sensors are run simultaneously while monitoring bus stability, sampling consistency, CPU utilization, memory usage, and timing accuracy.

\begin{checklist}
    \item Continuous operation without crashes.
    \item Stable timestamps.
    \item No sensor starvation.
\end{checklist}

\subsection{Stage 4 --- Integrated Audio Module}
\label{sec:stage4-audio-integration}

Verification covers continuous microphone capture, ambient light acquisition, and concurrent wireless communication.

\begin{checklist}
    \item No dropped audio buffers.
    \item Stable light sensor readings.
\end{checklist}

\subsection{Stage 5 --- Full System Integration}
\label{sec:stage5-full-integration}

Both modules operate together, with verification of time synchronization, data transmission, backend reception, and session management.

\begin{checklist}
    \item Complete synchronized datasets.
    \item No packet corruption.
\end{checklist}

\subsection{Functional Testing}
\label{sec:functional-testing}

Functional testing confirms that the system can:
\begin{checklist}
    \item Boot successfully.
    \item Detect all sensors.
    \item Acquire data.
    \item Timestamp correctly.
    \item Transmit data.
    \item Recover from temporary communication loss.
    \item Shut down safely.
\end{checklist}

\subsection{Stress Testing}
\label{sec:stress-testing}

\textbf{Tests conducted.}
\begin{itemize}
    \item 1-hour continuous operation.
    \item Extended operation (target duration based on battery capacity).
    \item Rapid reboot cycles.
    \item Repeated connection/disconnection tests.
    \item Continuous wireless transmission.
\end{itemize}

\textbf{Monitored quantities.}
\begin{itemize}
    \item Temperature.
    \item Memory usage.
    \item Packet loss.
    \item Unexpected resets.
\end{itemize}

\subsection{Power Validation}
\label{sec:power-validation}

\textbf{Verification.}
\begin{itemize}
    \item Battery charging.
    \item Battery discharge.
    \item Low-battery detection.
    \item Brownout handling.
    \item Controlled shutdown.
\end{itemize}

\textbf{Recorded quantities.}
\begin{itemize}
    \item Battery voltage.
    \item Runtime.
    \item Charging time.
    \item Device temperature.
\end{itemize}

\subsection{Sensor Calibration}
\label{sec:sensor-calibration}

Calibration procedures are documented for MAX30102, GSR, MCP9808, TEMT6000, and MPU6050. Each calibration record should include:
\begin{itemize}
    \item Calibration date.
    \item Firmware version.
    \item Calibration method.
    \item Environmental conditions.
\end{itemize}

\subsection{Software Validation}
\label{sec:software-validation}

\begin{checklist}
    \item Firmware boots correctly.
    \item Sensor drivers initialize.
    \item Communication stack functions.
    \item Logging system operates.
    \item Error recovery works.
\end{checklist}

\subsection{AI Validation}
\label{sec:ai-validation-ch8}

\begin{checklist}
    \item Feature extraction.
    \item Timestamp synchronization.
    \item Inference pipeline.
    \item Model confidence reporting.
    \item Handling of missing or corrupted data.
\end{checklist}

\subsection{Documentation Package}
\label{sec:documentation-package}

The complete documentation package should include:
\begin{checklist}
    \item Bill of Materials (BOM).
    \item Wiring diagrams.
    \item GPIO map.
    \item Power tree.
    \item Firmware architecture.
    \item AI architecture.
    \item Testing reports.
    \item Calibration records.
    \item Risk assessment.
    \item User manual.
    \item Assembly guide.
    \item Version history.
\end{checklist}

\subsection{Competition Readiness Checklist}
\label{sec:competition-readiness}

\begin{requirementbox}
\textbf{Before submission:}
\begin{checklist}
    \item Hardware fully assembled.
    \item Firmware stable.
    \item AI pipeline validated.
    \item Documentation complete.
    \item Safety review completed.
    \item Demonstration rehearsed.
    \item Backup firmware available.
    \item Backup hardware available (where practical).
\end{checklist}
\end{requirementbox}

\subsection{Known Remaining Risks}
\label{sec:known-remaining-risks-ch8}

\begin{observationbox}
\textbf{Current engineering risks include:}
\begin{itemize}
    \item GSR analog front-end refinement.
    \item Battery runtime limitations.
    \item Wireless reliability in congested environments.
    \item Motion artifacts affecting physiological signals.
    \item Audio quality in noisy environments.
    \item Long-term sensor stability.
\end{itemize}
Each risk should have a documented mitigation strategy.
\end{observationbox}

\subsection{Final Deliverables}
\label{sec:final-deliverables}

The complete PSYCON package should contain:
\begin{itemize}
    \item Hardware prototype.
    \item Embedded firmware.
    \item AI models.
    \item Backend software.
    \item Mobile/web interface (if applicable).
    \item Technical documentation.
    \item Validation and testing reports.
    \item User and maintenance guides.
    \item Safety and ethics documentation.
    \item Research report.
\end{itemize}

\subsection{Volume 1 Summary}
\label{sec:volume1-summary}

\begin{validationbox}
\textbf{This concludes the core engineering design document, covering:}
\begin{itemize}
    \item Project overview
    \item Hardware architecture
    \item Electrical design
    \item Firmware architecture
    \item AI pipeline
    \item System integration
    \item Verification and validation
\end{itemize}
\end{validationbox}
